% chapterintro
% ==============================
% 第6章 导言：chapter6-0-chinese.md
上一章，我们了解了基础STL算法，并且使用简单的例子实践操作了一下这些STL接口：大多数STL算法都将一个或多个迭代器对，作为其输入或输出参数。STL算法也能接受谓词函数，自定义比较函数和转换函数。最后，有些接口会返回迭代器，因为其他算法还会用到这个迭代器。

STL算法旨在保持简单和通用。这样大多数代码就可以使用STL算法，从而让代码看起来简单明了。一个经验丰富的C++开发者对于STL算法非常了解，其会在代码中尽可能使用STL算法，这会让其他人更容易明白这段代码存在的原因，这样能帮助开发者和阅读者产生最大程度的共鸣。一个开发者的大脑会很容易的解析这些普及度很高的算法，这要比解析一段复杂的循环简单的多。虽然实现的主要方式一样，但是细节方面还有不同。

我们使用STL数据结构结构能够很好的避免指针、裸数组和粗犷的结构体。那么接下来我们将升级对STL算法的理解，以便使用通用的STL算法来替代复杂的循环控制符合代码块。因为代码变得短小，在提高可读性的同时，又增加了通用性。这就能避免循环，仅调用\texttt{std}命名空间的算法即可，不过有时也会造成很糟糕的代码。我们不会去衡量代码是否糟糕，只会讨论可能性。

本章，我们将使用STL算法，以创造性的视野去了解现代C++能做些什么。我们将会实现属于我们自己的类STL算法，其能和已存在的数据结构完美结合，并且其他算法也会以同样的方式进行设计。我们也会将现有的STL算法与新算法相融合。这样的结合可以塑造出更加复杂的算法，不过其实现会更加短小，更具有可读性。过程中，我们可以看到STL是多么的简单和优雅。只有了解了所有方法，才能在使用时选择最合适的那一种。


% --- recipe 1: chapter6-1-chinese.md ---
\section{使用STL算法实现单词查找树类}

所谓的trie数据类型，能够对感兴趣的数据进行存储，并且易于查找。将文本语句分割成多个单词，放置在列表中，我们能发现其开头一些单词的共性。

让我们看下下图，这里有两个句子“hi how are you”和“hi how do you do”，存储在一个类似于树的结构体中。其都以“hi how”开头，句子后面不同的部分，划分为树结构：

\myGraphic{0.9}{images/chapter6/6-1-1.png}{}

因为trie数据结构结合了相同的前缀，其也称为前缀树，很容易使用STL的数据结构实现。本章我们将关注如何实现我们自己的trie类。

\mySmallsection{How to do it...}

本节，我们将使用STL数据结构和算法实现前缀树结构。

\begin{enumerate}
\item 包含必要的头文件和声明所使用的命名空间

\begin{codebox}[cpp]
#include <iostream>
#include <optional>
#include <algorithm>
#include <functional>
#include <iterator>
#include <map>
#include <vector>
#include <string>
using namespace std;
\end{codebox}

\item 我们首先实现一个类。我们的实现中，trie为\texttt{map}的递归映射。每个trie节点够包含一个\texttt{map}，节点的有效值\texttt{T}映射了下一个节点：

\begin{codebox}[cpp]
template <typename T>
class trie
{
    map<T, trie> tries;
\end{codebox}

\item 将新节点插入队列的代码很简单。使用者需要提供一个\texttt{begin/end}迭代器对，并且会通过循环进行递归。当用户输入的序列为{1, 2, 3}时，我们可以将1作为一个子trie，2为下一个子trie，以此类推。如果这些子trie在之前不存在，其将会通过\texttt{std::map}的\texttt{[]}操作符进行添加：

\begin{codebox}[cpp]
public:
    template <typename It>
    void insert(It it, It end_it) {
        if (it == end_it) { return; }
        tries[*it].insert(next(it), end_it);
    }
\end{codebox}

\item 我们这里也会定义一个辅助函数，用户只需要提供一个容器，之后辅助函数就会通过迭代器自动进行查询：

\begin{codebox}[cpp]
    template <typename C>
    void insert(const C &container) {
        insert(begin(container), end(container));
    } 
\end{codebox}

\item 调用我们的类时，可以写成这样\texttt{my\_trie.insert(\{"a", "b","c"\}); }，必须帮助编译器正确的判断出这段代码中的所有类型，因此我们又添加了一个函数，这个函数用于重载的插入接口：

\begin{codebox}[cpp]
    void insert(const initializer_list<T> &il) {
        insert(begin(il), end(il));
    }
\end{codebox}

\item 我们也想了解，trie中有什么，所以我们需要一个打印函数。为了打印，我们可以对tire进行深度遍历。这样根节点下面的是第一个叶子节点，我们会记录我们所看到的元素的负载。当我们达到叶子节点，那么就可以进行打印了。我们会看到，当到达叶子的时候\texttt{tries.empty()}为true。递归调用print后，我们将再次弹出最后添加的负载元素：

\begin{codebox}[cpp]
void print(vector<T> &v) const {
    if (tries.empty()) {
        copy(begin(v), end(v),
            ostream_iterator<T>{cout, " "});
        cout << '\n';
    }
    for (const auto &p : tries) {
        v.push_back(p.first);
        p.second.print(v);
        v.pop_back();
    }
}
\end{codebox}

\item 打印函数需要传入一个可打印负载元素的列表，不过用户不需要传入任何参数就能调用它。这样，我们就定义了一个无参数的打印函数，其构造了辅助列表对象：

\begin{codebox}[cpp]
    void print() const {
        vector<T> v;
        print(v);
    } 
\end{codebox}

\item 现在，我们就可以创建和打印trie了，我们将先搜索子trie。当trie包含的序列为\texttt{\{a, b, c\}}和\texttt{\{a, b, d, e\}}，并且我们给定的序列为\texttt{\{a, b\}}，对于查询来说，返回的子序列为包含\texttt{\{c\}}和\texttt{\{d, e\}}的部分。当我们找到子trie，将返回一个\texttt{const}的引用。在搜索中，也会出现没有要搜索序列的情况。即便如此，我们还是要返回些什么。\texttt{std::optional}是一个非常好的帮手，因为当没有找到匹配的序列，我们可以返回一个空的\texttt{optional}对象：

\begin{codebox}[cpp]
template <typename It>
optional<reference_wrapper<const trie>>
subtrie(It it, It end_it) const {
    if (it == end_it) { return ref(*this); }
    auto found (tries.find(*it));
    if (found == end(tries)) { return {}; }
    return found->second.subtrie(next(it), end_it);
}
\end{codebox}

\item 与\texttt{insert}方法类似，我们将提供一个只需要一个参数的\texttt{subtrie}方法，其能自动的从输入容器中获取迭代器：

\begin{codebox}[cpp]
    template <typename C>
    auto subtrie(const C &c) {
        return subtrie(begin(c), end(c));
    }
};
\end{codebox}

\item 这样就实现完了。我们在主函数中使用我们trie类，使用\texttt{std::string}类型对类进行特化，并实例化对象：

\begin{codebox}[cpp]
int main()
{
    trie<string> t;
    t.insert({"hi", "how", "are", "you"});
    t.insert({"hi", "i", "am", "great", "thanks"});
    t.insert({"what", "are", "you", "doing"});
    t.insert({"i", "am", "watching", "a", "movie"});
\end{codebox}

\item 打印整个trie：

\begin{codebox}[cpp]
    cout << "recorded sentences:\n";
    t.print();
\end{codebox}

\item 而后，我们将获取输入语句的子trie，其以“hi”开头：

\begin{codebox}[cpp]
    cout << "\npossible suggestions after \"hi\":\n";
    if (auto st (t.subtrie(initializer_list<string>{"hi"}));
        st) {
        st->get().print();
    }
}
\end{codebox}

\item 编译并运行程序，其会返回两个句子的以“hi”开头的子trie：

\begin{codebox}[cpp]
$ ./trie
recorded sentences:
hi how are you
hi i am great thanks
i am watching a movie
what are you doing
possible suggestions after "hi":
how are you
i am great thanks
\end{codebox}

\end{enumerate}

\mySmallsection{How it works...}

有趣的是，单词序列的插入代码要比在子trie查找给定字母序列的代码简单许多。所以，我们首先来看一下插入代码：

\begin{codebox}[cpp]
template <typename It>
void trie::insert(It it, It end_it) {
    if (it == end_it) { return; }
    tries[*it].insert(next(it), end_it);
}
\end{codebox}

迭代器对\texttt{it}和\texttt{end\_it}，表示要插入的字符序列。\texttt{tries[*it]}代表在子trie中要搜索的第一个字母，然后调用\texttt{.insert(next(it), end\_it);}对更低级的子trie序列使用插入函数，使用迭代器一个词一个词的推进。\texttt{ if (it == end\_it) \{ return; \}}行会终止递归。返回语句不会做任何事情，这到有点奇怪了。所有插入操作都在\texttt{tries[*it]}语句上进行，\texttt{std::map}的中括号操作将返回键所对应的值或是创建该键，相关的值(本节中映射类型是一个trie)由默认构造函数构造。这样，当我们查找不理解的单词时，就能隐式的创建一个新的trie分支。

查找子trie看起来十分复杂，因为我们没有必要隐藏那么多的代码：

\begin{codebox}[cpp]
template <typename It>
optional<reference_wrapper<const trie>>
subtrie(It it, It end_it) const {
    if (it == end_it) { return ref(*this); }
    auto found (tries.find(*it));
    if (found == end(tries)) { return {}; }

    return found->second.subtrie(next(it), end_it);
}
\end{codebox}

这段代码的主要部分在于\texttt{auto found (tries.find(*it));}。我们使用find来替代中括号操作符。当我们使用中括号操作符进行查找时，trie将会为我们创建丢失的元素(顺带一提，当我们尝试这样做，类的函数为\texttt{const}，所以这样做事不可能的。这样的修饰能帮助我们减少bug的发生)。

另一个细节是返回值\texttt{optional<reference\_wrapper<const trie>>}。我们选择\texttt{std::optional}作为包装器，因为其可能没有我们所要找打tire。当我们仅插入“hello my friend”，那么就不会找到“goodbye my friend”。这样，我们仅返回\texttt{\{\}}就可以了，其代表返回一个空\texttt{optional}对象给调用者。不过这还是没有解释，我们为什么使用\texttt{reference\_wrapper}代替\texttt{optional<const trie \&>}。\texttt{optional}的实例，其为\texttt{trie\&}类型，是不可赋值的，因此不会被编译。使用\texttt{reference\_warpper}实现一个引用，就是用来对对象进行赋值的。


% --- recipe 2: chapter6-2-chinese.md ---
\section{使用树实现搜索输入建议生成器}

上网时，在搜索引擎中输入要查找的东西时，对应下拉选项中会尝试猜测你想要查找什么。这种猜测是基于之前相关主题被查找的数量。有时搜索引擎十分有趣，其会显示一些奇怪的主题。

\myGraphic{0.9}{images/chapter6/6-2-1.png}{}

本章，我们将使用树类实现一个简单的搜索建议引擎。

\mySmallsection{How to do it...}

本节，我们将实现一个终端应用，其能接受输入，并且能对所要查找的内容进行猜测，当然猜测的依据是我们用文本完成的“数据库”。

\begin{enumerate}
\item 包含必要的头文件和声明所使用的命名空间：

\begin{codebox}[cpp]
#include <iostream>
#include <optional>
#include <algorithm>
#include <functional>
#include <iterator>
#include <map>
#include <list>
#include <string>
#include <sstream>
#include <fstream>
using namespace std;
\end{codebox}

\item 我们将使用上一节实现的trie类：

\begin{codebox}[cpp]
template <typename T>
class trie
{
    map<T, trie> tries;
public:
    template <typename It>
    void insert(It it, It end_it) {
        if (it == end_it) { return; }
        tries[*it].insert(next(it), end_it);
    }
    template <typename C>
    void insert(const C &container) {
        insert(begin(container), end(container));
    }
    void insert(const initializer_list<T> &il) {
        insert(begin(il), end(il));
    }
    void print(list<T> &l) const {
        if (tries.empty()) {
            copy(begin(l), end(l),
                ostream_iterator<T>{cout, " "});
            cout << '\n';
        }
        for (const auto &p : tries) {
            l.push_back(p.first);
            p.second.print(l);
            l.pop_back();
        }
    }
    void print() const {
        list<T> l;
        print(l);
    }
    template <typename It>
    optional<reference_wrapper<const trie>>
    subtrie(It it, It end_it) const {
        if (it == end_it) { return ref(*this); }
        auto found (tries.find(*it));
        if (found == end(tries)) { return {}; }
        return found->second.subtrie(next(it), end_it);
    }
    template <typename C>
    auto subtrie(const C &c) const {
        return subtrie(begin(c), end(c));
    }
};
\end{codebox}

\item 实现一个简单的辅助函数，这个函数将用于提示用户输入他们想要查找的东西：

\begin{codebox}[cpp]
static void prompt()
{
    cout << "Next input please:\n > ";
} 
\end{codebox}

\item 主函数中，我们打开一个文本文件，其作为我们的基础数据库。我们逐行读取文本文件的内容，并且将数据放入trie中解析：

\begin{codebox}[cpp]
int main()
{
    trie<string> t;
    fstream infile {"db.txt"};
    for (string line; getline(infile, line);) {
        istringstream iss {line};
        t.insert(istream_iterator<string>{iss}, {});
    }
\end{codebox}

\item 现在可以使用构建好的trie类，并且需要实现接收用户查询输入的接口。会提示用户进行输入，并且将用户的输入整行读取：

\begin{codebox}[cpp]
prompt();
for (string line; getline(cin, line);) {
    istringstream iss {line};
\end{codebox}

\item 通过文本输入，可以使用trie对其子trie进行查询。如果在数据库中已经有相应的语句，那么会对输入进行建议，否则会告诉用户没有建议给他们：

\begin{codebox}[cpp]
if (auto st (t.subtrie(istream_iterator<string>{iss}, {}));
    st) {
    cout << "Suggestions:\n";
    st->get().print();
} else {
    cout << "No suggestions found.\n";
}
\end{codebox}

\item 之后，将打印一段分割符，并且再次等待用户的输入：

\begin{codebox}[cpp]
        cout << "----------------\n";
        prompt();
    }
}
\end{codebox}

\item 运行程序之前，我们需要将db.txt文件进行设置。查找的输入可以是任何字符，并且其不确保是已经排过序的。进入trie类的所有语句：

\begin{codebox}[cpp]
do ghosts exist
do goldfish sleep
do guinea pigs bite
how wrong can you be
how could trump become president
how could this happen to me
how did bruce lee die
how did you learn c++
what would aliens look like
what would macgiver do
what would bjarne stroustrup do
...
\end{codebox}

\item 创建完db.txt之后，我们就可以运行程序了。其内容如下所示：

\begin{codebox}[cpp]
hi how are you
hi i am great thanks
do ghosts exist
do goldfish sleep
do guinea pigs bite
how wrong can you be
how could trump become president
how could this happen to me
how did bruce lee die
how did you learn c++
what would aliens look like
what would macgiver do
what would bjarne stroustrup do
what would chuck norris do
why do cats like boxes
why does it rain
why is the sky blue
why do cats hate water
why do cats hate dogs
why is c++ so hard
\end{codebox}

\item 编译并运行程序，然后进行输入查找：

\begin{codebox}[cpp]
$ ./word_suggestion
Next input please:
> what would
Suggestions:
aliens look like
bjarne stroustrup do
chuck norris do
macgiver do
----------------
Next input please:
> why do
Suggestions:
cats hate dogs
cats hate water
cats like boxes
----------------
Next input please:
>
\end{codebox}

\mySmallsection{How it works...}

\end{enumerate}

trie是如何工作的，已经在上一节中介绍过了，不过本节我们对其进行填充和查找的过程看起来有些奇怪。让我们来仔细观察一下代码片段，其使用文本数据库文件对空trie类进行填充：

\begin{codebox}[cpp]
fstream infile {"db.txt"};
for (string line; getline(infile, line);) {
    istringstream iss {line};
    t.insert(istream_iterator<string>{iss}, {});
}
\end{codebox}

这段代码会逐行的将文本文件中的内容读取出来。然后，我们将字符串拷贝到一个\texttt{istringstream}对象中。我们可以根据输入流对象，创建一个\texttt{istring\_iterator}迭代器，其能帮助我们查找子trie。这样，我们就不需要将字符串放入\texttt{vector}或\texttt{list}中了。上述代码中，有一段不必要的内存分配，可以使用移动方式，将\texttt{line}中的内容移动到iss中，避免不必要的内存分配。不过，\texttt{std::istringstream}没有提供构造函数，所以只能将\texttt{std::string}中的内容移动到流中。不过，这里会对输入字符串进行复制。

当在trie中查询用户的输入时，使用了相同的策略，但不使用输入文件流。我们使用\texttt{std::cin}作为替代，因为\texttt{trie::subtrie}对迭代器的操作，和\texttt{trie::insert}如出一辙。

\mySmallsection{There's more...}

这里有必要对每个trie节点添加统计变量，这样我们就能知道各种前缀被查询的频率。因此，我们就可以将程序的建议进行排序，当前的搜索引擎就是这样做的。智能手机触摸屏文本输入的建议，也可以通过这种方式实现。

这个修改就留给读者当作业了。 ：）


% --- recipe 3: chapter6-3-chinese.md ---
\section{使用STL数值算法实现傅里叶变换}

信号处理领域傅里叶变换是非常重要和著名的公式。这个公式发现于200年前，其计算机用例实很多了。傅里叶变换可以用于音频/图像/视频压缩、音频滤波、医疗图像设备和用于辨识声音的手机引用。

因为其应用领域广泛，STL也试图将其用在数值计算领域。傅里叶变换只是其中一个例子，同样也是非常棘手的一个。其公式如下所示：

\myGraphic{0.9}{images/chapter6/6-3-1.png}{}

公式基于累加和的变换。累加中的每个元素是输入信号向量中的一个数据点和表达式\texttt{exp(-2 * i * ...)}的乘积。这里需要一些工程数学的知识，你需要简单的了解复数的概念，如果你没有相关的知识，了解概念就可以了。仔细观察这个公式，其就是将信号中的所有数据点进行加和(信号数据的长度为N)，其循环索引值为j。其中k是另一个循环变量，因为傅里叶变换计算出的是一组值。在这组值中，每一个数据点都表示着一段重复波形的幅值和相位，这些信息不包含在原始数据中。当使用循环对其进行实现时，代码可能就会写成下面这样：

\begin{codebox}[cpp]
csignal fourier_transform(const csignal &s) {
    csignal t(s.size());
    const double pol {-2.0 * M_PI / s.size()};
    for (size_t k {0}; k < s.size(); ++k) {
        for (size_t j {0}; j < s.size(); ++j) {
            t[k] += s[j] * polar(1.0, pol * k * j);
        }
    }
    return t;
}
\end{codebox}

这里\texttt{csignal}的类型可能是\texttt{std::vector}，其每个元素都是一个复数。对于复数而言，STL中已经有了对应的数据结构可以对其进行表示——\texttt{std::complex}。\texttt{std::polar}函数计算得是\texttt{exp(-2 * i * ...)}部分。

这样实现看起来也挺好，不过本节中我们将使用STL工具对其进行实现。

\mySmallsection{How to do it...}

本节，我们将实现傅里叶变换和逆变换，然后会对一些信号进行转换：

\begin{enumerate}
\item 首先，包含必要的头文件和声明所使用的命名空间：

\begin{codebox}[cpp]
#include <iostream>
#include <complex>
#include <vector>
#include <algorithm>
#include <iterator>
#include <numeric>
#include <valarray>
#include <cmath>
using namespace std;
\end{codebox}

\item 信号点的值一个复数，我们使用\texttt{std::complex}来表示，并使用\texttt{double}进行特化。我们可以对类型进行别名操作，使用\texttt{cmple}表示两个\texttt{double}值，这两个\texttt{double}值分别表示复数的实部和虚部。使用\texttt{csdignal}来别名相应的\texttt{vector}对象：

\begin{codebox}[cpp]
using cmplx = complex<double>;
using csignal = vector<cmplx>;
\end{codebox}

\item 我们需要使用数值指针遍历数值序列。公式中的变量k和j就会随着序列进行累加：

\begin{codebox}[cpp]
class num_iterator {
    size_t i;
public:
    explicit num_iterator(size_t position) : i{position} {}
    size_t operator*() const { return i; }
    num_iterator& operator++() {
        ++i;
        return *this;
    }
    bool operator!=(const num_iterator &other) const {
        return i != other.i;
    }
};
\end{codebox}

\item 傅里叶变换需要接收一个信号，并返回一个新的信号。返回的信号表示已经经过傅里叶变换的信号。通过傅里叶逆变换，我们可以将一个经过傅里叶变换的信号，还原成原始信号，这里我们会提供一个可选的\texttt{bool}参数，其会决定变换的方向。\texttt{bool}参数作为参数是一种不好习惯，特别是在一个函数的签名中出现多次。我们这有个很简洁的例子。我们做的第一件事，是使用原始信号的尺寸来分配新的信号数组：

\begin{codebox}[cpp]
csignal fourier_transform(const csignal &s, bool back = false)
{
    csignal t (s.size());
\end{codebox}

\item 公式中有两个因子，其看起来是相同的。让我们将其打包成一个变量：

\begin{codebox}[cpp]
    const double pol {2.0 * M_PI * (back ? -1.0 : 1.0)};
    const double div {back ? 1.0 : double(s.size())};
\end{codebox}

\item \texttt{std::accumulate}很适合用来执行公式中的累加部分，我们将对一个范围内的数值使用\texttt{accumulate}。对于每个值，我们将逐步的进行单个相加。\texttt{std::accumulate}算法会调用一个二元函数。该函数的第一个参数为目前为止我们所累加的变量\texttt{sum}，第二个参数为范围内下一个要累加的值。我们会在信号\texttt{s}中对当前为止的值进行查找，并且会将其和复数因子\texttt{pol}相乘。然后，我们返回新的\texttt{sum}。这里的二元函数，使用Lambda表达式进行包装，因为我们将在每次\texttt{accumulate}的调用时，\texttt{j}变量的值是不同的。因为其是二维循环算法，所以内层Lambda做内部的循环，外层Lambda做外层的循环：

\begin{codebox}[cpp]
auto sum_up ([=, &s] (size_t j) {
    return [=, &s] (cmplx c, size_t k) {
        return c + s[k] *
            polar(1.0, pol * k * j / double(s.size()));
    };
});
\end{codebox}

\item 傅里叶的内部循环，现在使用\texttt{std::accumulate}进行，算法中每个位置都会进行加和。我们使用Lambda表达式来实现，这样我们就能计算出傅里叶变换数组中的每个数据点的值：

\begin{codebox}[cpp]
    auto to_ft ([=, &s](size_t j){
        return accumulate(num_iterator{0},
                          num_iterator{s.size()},
                          cmplx{},
                          sum_up(j))
        / div;
    });
\end{codebox}

\item 目前位置，还没有执行傅里叶变换的代码。我们会准备大量的功能性代码，他们会帮助我们完成很多事情。\texttt{std::transform}的调用将会使j的值在[0, N)间变换(这步是在外层循环完成)。变换之后的值将全部放入\texttt{t}中，\texttt{t}就是我们要返回给用户的值：

\begin{codebox}[cpp]
    transform(num_iterator{0}, num_iterator{s.size()},
             begin(t), to_ft);
    return t;
}
\end{codebox}

\item 我们将会实现一些辅助函数帮助我们生成信号。首先实现的是一个余弦信号生成器，其会返回一个Lambda表达式，这个表达式通过传入的长度参数，产生对应长度的余弦信号数据。信号本身的长度是不固定的，但是其有固定的周期。周期为N，意味着该信号会在N步之后重复。返回的Lambda表达式不接受任何参数。我们可以重复的对其进行调用，并且每次调用表达式将会返回给我们下一个时间点的信号值：

\begin{codebox}[cpp]
static auto gen_cosine (size_t period_len){
    return [period_len, n{0}] () mutable {
        return cos(double(n++) * 2.0 * M_PI / period_len);
    };
}
\end{codebox}

\item 我们所要生成另一个波形是方波。该波形会在\texttt{-1}和\texttt{+1}两值间震荡，其中不会有其他的值。公式看起来有点复杂，但是其变换非常简单，也就是将值n置为\texttt{+1}或\texttt{-1}，并且其震荡周期为\texttt{period\_len}。这里要注意，我们没有使用0对n进行初始化。这样，我们的方波的其实位置就在\texttt{+1}上：

\begin{codebox}[cpp]
static auto gen_square_wave (size_t period_len)
{
    return [period_len, n{period_len*7/4}] () mutable {
        return ((n++ * 2 / period_len) % 2) * 2 - 1.0;
    };
}
\end{codebox}

\item 产生实际信号可以通过\texttt{vector}和信号生成器联合进行，使用重复调用信号生成器对\texttt{vector}数组进行填充。\texttt{std::generate}就用来完成这个任务的。其接受一组\texttt{begin/end}迭代器组和一个生成函数。对于每个合法的迭代器，都会进行\texttt{*it = gen()}。通过将这些代码包装成一个函数，我们可以很容易的生成一个信号数组：

\begin{codebox}[cpp]
template <typename F>
static csignal signal_from_generator(size_t len, F gen)
{
    csignal r (len);
    generate(begin(r), end(r), gen);
    return r;
}
\end{codebox}

\item 最后，我们需要将信号的结果进行打印。我们可以将数组中的值拷贝到输出流迭代器中进行输出，不过我们需要先将数据进行变换，因为我们的信号数据都是复数对。这样，我们只需要在意每个点的实部就好；所以，我们可以将数组扔到\texttt{std::transform}中进行变换，然后将实部提取出来：

\begin{codebox}[cpp]
static void print_signal (const csignal &s)
{
    auto real_val ([](cmplx c) { return c.real(); });
    transform(begin(s), end(s),
             ostream_iterator<double>{cout, " "}, real_val);
    cout << '\n';
}
\end{codebox}

\item 目前为止，傅里叶公式就已经实现了，不过现在还没有信号进行变换。这个工作我们将在主函数中完成。我们先来定义信号数据的长度：

\begin{codebox}[cpp]
int main()
{
    const size_t sig_len {100};
\end{codebox}

\item 现在来生成信号数据，转换他们，然后进行打印。首先，生成一个余弦信号和一个方波信号。这两组信号的长度和周期数相同：

\begin{codebox}[cpp]
auto cosine (signal_from_generator(sig_len,
    gen_cosine( sig_len / 2)));
auto square_wave (signal_from_generator(sig_len,
    gen_square_wave(sig_len / 2)));
\end{codebox}

\item 那么现在有了两个波形信号。为了生成第三个信号，我们对方波信号进行傅里叶变换，并且保存在\texttt{trans\_sqw}数组中。方波的傅里叶变换有些特殊，我们在后面会进行介绍。索引从10到\texttt{(signal\_length - 10)}都设置为0.0。经过傅里叶变换之后，原始信号将发生很大的变化。我们将在最后看到结果：

\begin{codebox}[cpp]
auto trans_sqw (fourier_transform(square_wave));
fill (next(begin(trans_sqw), 10), prev(end(trans_sqw), 10), 0);
auto mid (fourier_transform(trans_sqw, true));
\end{codebox}

\item 现在，我们有三个信号：余弦、mid和方波。对于每个信号，我们将会打印其原始波形，和傅里叶变换过后的波形。输出将有六条曲线组成：

\begin{codebox}[cpp]
    print_signal(cosine);
    print_signal(fourier_transform(cosine));
    print_signal(mid);
    print_signal(trans_sqw);
    print_signal(square_wave);
    print_signal(fourier_transform(square_wave));
}
\end{codebox}

\item 编译并运行程序，终端上会打印出大量的数据。如果这里使用绘图输出，就可以看到如下的结果：

\myGraphic{0.9}{images/chapter6/6-3-2.png}{}

\end{enumerate}

\mySmallsection{How it works...}

这段代码又两个比较难理解的部分。第一个是傅里叶变换本身，另一个是使用可变Lambda表达式生成信号数据。

首先，我们来看一下傅里叶变换。其核心部分在循环中实现(虽然没有在我们实现中这样做，但可以结合代码看下介绍中的公式)，可能会以如下方式实现：

\begin{codebox}[cpp]
for (size_t k {0}; k < s.size(); ++k) {
    for (size_t j {0}; j < s.size(); ++j) {
        t[k] += s[j] * polar(1.0, pol * k * j / double(s.size()));
    }
}
\end{codebox}

基于STL算法\texttt{std::transform}和\texttt{std::accumulate}，我们完成了自己的例子，总结一下就类似如下的伪代码：

\begin{codebox}[text]
transform(num_iterator{0}, num_iterator{s.size()}, ...
    accumulate((num_iterator0}, num_iterator{s.size()}, ...
        c + s[k] * polar(1.0, pol * k * j / double(s.size()));
\end{codebox}

和循环相比，结果完全一样。当然，使用STL算法也可以产生不太好的代码。不管怎么样吧，这个实现是不依赖所选用的数据结构。其对于列表也起作用(虽然这没有太大的意义)。另一个好处是，在C++17中STL很容易并行(将在本书的另一个章节进行介绍)，当需要并行的时候，我们就需要对纯循环进行重构和拆分，将其放入指定的线程中(除非使用类似OpenMP这样的并发库，其会自动的将循环进行重构)。

下一个难点是信号生成。让我来看一下另一个\texttt{gen\_cosine}:

\begin{codebox}[cpp]
static auto gen_cosine (size_t period_len)
{
    return [period_len, n{0}] () mutable {
        return cos(double(n++) * 2.0 * M_PI / period_len);
    };
}
\end{codebox}

每一个Lambda表达式代表一个函数对象，其会在每次调用时改变自身的状态。其状态包含两个变量\texttt{period\_len}和\texttt{n}。变量n会在每次调用时，进行变更。在不同的时间点上，得到的是不同的信号值，并且在时间增加时会使用\texttt{n++}对\texttt{n}的值进行更新。为了获得信号值的数组，我们创建了辅助函数\texttt{signal\_from\_generator}：

\begin{codebox}[cpp]
template <typename F>
static auto signal_from_generator(size_t len, F gen)
{
    csignal r (len);
    generate(begin(r), end(r), gen);
    return r;
}
\end{codebox}

这个函数会通过所选长度创建一个信号\texttt{vector}，并且调用\texttt{std::generate}对数据点进行填充。数组r中的每一个元素，都会调用一个\texttt{gen}函数。\texttt{gen}函数是是一种自修改函数对象，我们使用相同的方式创建了\texttt{gen\_cosine}对象。

\begin{myInfoBox}
Note： 本节例子中，STL没有让代码更加的优雅。如果将范围库添加入STL(希望在C++20时加入)，那么可能就会有改观。
\end{myInfoBox}


% --- recipe 4: chapter6-4-chinese.md ---
\section{计算两个vector的误差和}

对两个值进行计算的时候，计算机的计算结果与我们期望的结果有一定的差别。比如，测量由多个数据点组成的信号之间的差异，通常会涉及相应数据点的循环和减法等计算。

我们给出一个简单的计算信号a与信号b之间的误差公式：

\myGraphic{0.9}{images/chapter6/6-4-1.png}{}

对于每一个\texttt{i}，都会计算一次\texttt{a[i] - b[i]}，对差值求平方(负值和正值就能进行比较)，最后计算平方差的和。通常我们会使用循环来做这件事，但是为了让事情更加好玩，我们决定使用STL算法来完成。使用STL的好处是，无需耦合特定的数据结果。我们的算法能够适应\texttt{vector}和类似链表的数据结构，不用直接进行索引。

\mySmallsection{How to do it...}

本节，我们将创建两个信号，并计算这两个信号之间的误差：

\begin{enumerate}
\item 依旧是包含必要的头文件和声明所使用的命名空间。

\begin{codebox}[cpp]
#include <iostream>
#include <cmath>
#include <algorithm>
#include <numeric>
#include <vector>
#include <iterator>
using namespace std; 
\end{codebox}

\item 我们将对两个信号的误差和进行计算。这两个信号一个是\texttt{sine}，另一个信号也是\texttt{sine}，不过其中之一的使用\texttt{double}类型进行保存，另一个使用\texttt{int}类型进行保存。因为\texttt{double}和\texttt{int}类型表示数值的范围有差异，就像是模拟信号\texttt{as}转换成数字信号\texttt{ds}。

\begin{codebox}[cpp]
int main()
{
    const size_t sig_len {100};
    vector<double> as (sig_len); // a for analog
    vector<int> ds (sig_len); // d for digital
\end{codebox}

\item 为了生成一个\texttt{sin}波形，我们事先了一个简单的Lambda表达式，并可以传入一个可变的计数变量\texttt{n}。我们可以经常在需要的时候调用表达式，其将返回下一个时间点的\texttt{sine}波形。\texttt{std::generate}可以使用信号值来填充数组，并且使用\texttt{std::copy}将数组中的\texttt{double}类型的变量，转换成\texttt{int}类型变量：

\begin{codebox}[cpp]
    auto sin_gen ([n{0}] () mutable {
        return 5.0 * sin(n++ * 2.0 * M_PI / 100);
    });
    generate(begin(as), end(as), sin_gen);
    copy(begin(as), end(as), begin(ds));
\end{codebox}

\item 我们可以对信号进行打印，也可以使用绘图进行显示：

\begin{codebox}[cpp]
    copy(begin(as), end(as),
        ostream_iterator<double>{cout, " "});
    cout << '\n';
    copy(begin(ds), end(ds),
        ostream_iterator<double>{cout, " "});
    cout << '\n'; 
\end{codebox}

\item 现在来计算误差和，我们使用\texttt{std::inner\_product}，因为这个函数能帮助我们计算两个信号矢量的差异。该函数能在指定范围内进行迭代，然后选择相应位置上进行差值计算，然后在进行平方，再进行相加：

\begin{codebox}[cpp]
    cout << inner_product(begin(as), end(as), begin(ds),
                        0.0, std::plus<double>{},
                        [](double a, double b) {
                            return pow(a - b, 2);
                        })
        << '\n';
}
\end{codebox}

\item 编译并运行程序，我们就能得到两条曲线，还有一条曲线代表的是两个信号的误差和。最终这两个信号的误差为40.889。当我们使用连续的方式对误差进行统计，要对值进行逐对匹配，然后得到无法曲线，其就像我们在下图中看到的一样：

\myGraphic{0.9}{images/chapter6/6-4-2.png}{}

\end{enumerate}

\mySmallsection{How it works...}

本节，我们需要将两个向量放入循环中，然后对不同位置的值计算差值，然后差值进行平方，最后使用\texttt{std::inner\_product}将差的平方进行加和。这样，我们可以使用Lambda表达式来完成求差值平方的操作——\texttt{[](double a, double b)\{return pow(a - b), 2\}}，这样就可以通过传入不同的参数来计算差值平方。

这里我们可以看下\texttt{std::inner\_product}是如何工作的：

\begin{codebox}[cpp]
template<class InIt1, class InIt2, class T, class F, class G>
T inner_product(InIt1 it1, InIt1 end1, InIt2 it2, T val,
               F bin_op1, G bin_op2)
{
    while(it1!= end1){
        val = bin_op1(val, bin_op2(*it1, *it2));
        ++it1;
        ++it2;
    }
    return value;
}
\end{codebox}

算法会接受一对\texttt{begin/end}迭代器作为第一个输入范围，另一个\texttt{begin}迭代器代表第二个输入范围。我们的例子中，这些迭代器所指向的是\texttt{vector}，并对这两个\texttt{vector}进行误差和的计算。\texttt{val}是一个初始化值。我们这里将其设置为\texttt{0.0}。然后，算法可以接受两个二元函数，分别为\texttt{bin\_op1}和\texttt{bin\_op2}。

我们会发现，这个算法与\texttt{std::accumulate}很相似。不过\texttt{std::accumulate}只对一个范围进行操作。当将\texttt{bin\_op2(*it1, *it2)}看做一个迭代器，那么我们可以简单的是用\texttt{accumulate}算法进行计算了。所以，我们可以将\texttt{std::inner\_product}看成是带有打包输入范围的\texttt{std::accumulate}。

例子中，打包函数就是\texttt{pow(a - b, 2)}。因为我们需要将所有元素的差平方进行加和，所以我们选择\texttt{std::plus<double>}作为\texttt{bin\_op1}。


% --- recipe 5: chapter6-5-chinese.md ---
\section{使用ASCII字符曼德尔布罗特集合}

1975年，数学家贝诺曼德尔布罗特(Benoit Mandelbrot)创造了一个术语——\textbf{分形}。分形是一个数学图像或者集合，这个术语中包含了很多有趣的数学特性，不过最后看起来分形更像是艺术品。分形图像看起来是无限重复的缩小。其中最为众人所知的分形是曼德尔布罗特(Mandelbrot)集合，其集合看起来就像下图一样：

\myGraphic{0.9}{images/chapter6/6-5-1.png}{}

曼德尔布罗特集合可以通过迭代下面的等式得到：

\myGraphic{0.9}{images/chapter6/6-5-2.png}{}

\texttt{z}和\texttt{c}变量都是复数。曼德尔布罗特集合包含等式所覆盖所有让方程收敛的\texttt{c}值，也就是海报彩色的部分。有些值收敛的早，有些值收敛的晚，这里用不同的颜色对这些值进行描述，所以我们能在海报中看到各种不同的颜色。对于那些不收敛的值，我们则直接将其所在区域直接涂黑。

使用STL的\texttt{std::complex}类，且不使用循环来实现上面的等式。这并不是炫技，只是为了让大家更容易理解STL相关特性的使用方式。

\mySmallsection{How to do it...}

本节，我们将打印类似墙上海报的图，不过是使用ASCII字符将图像打印在终端上：

\begin{enumerate}
\item 包含必要的头文件并声明所使用的命名空间：

\begin{codebox}[cpp]
#include <iostream>
#include <algorithm>
#include <iterator>
#include <complex>
#include <numeric>
#include <vector>
using namespace std;
\end{codebox}

\item 曼德尔布罗特集合和之前的等式，都是对复数进行操作。所以，我们需要使用类型别名，使用\texttt{cmplx}来代表\texttt{std::complex}，并特化为\texttt{double}类型：

\begin{codebox}[cpp]
using cmplx = complex<double>;
\end{codebox}

\item 我们将使用大约20行的代码来完成一个ASCII组成的曼德尔布罗特集合图像，不过我们会将逻辑逐步实现，最后将所有结果进行组合。第一步就是实现一个函数，用于将整型坐标缩放为浮点坐标。这也就是我们如何在屏幕上特定的位置上打印相应的字符。我们想要的是曼德尔布罗特集合中复数的坐标，就需要实现一个函数，用于将对应的坐标转换成相应的几何图形。用一个Lambda表达式来构建这些变量，并将其返回。该函数能将\texttt{int}类型的函数转换成一个\texttt{double}类型的函数：

\begin{codebox}[cpp]
static auto scaler(int min_from, int max_from,
double min_to, double max_to)
{
    const int w_from {max_from - min_from};
    const double w_to {max_to - min_to};
    const int mid_from {(max_from - min_from) / 2 + min_from};
    const double mid_to {(max_to - min_to) / 2.0 + min_to};
    return [=] (int from) {
        return double(from - mid_from) / w_from * w_to + mid_to;
    };
}
\end{codebox}

\item 现在需要在一个维度上进行坐标变换，不过曼德尔布罗特集合使用的是二维坐标系。为了能将(x, y)坐标系统转换成另一个，我们需要将\texttt{x-scaler}和\texttt{y-scaler}相结合，并且构建一个\texttt{cmplx}实例作为输出：

\begin{codebox}[cpp]
template <typename A, typename B>
static auto scaled_cmplx(A scaler_x, B scaler_y)
{
    return [=](int x, int y) {
        return cmplx{scaler_x(x), scaler_y(y)};
    };
}
\end{codebox}

\item 将坐标转换到正确的维度上后，就可以来实现曼德尔布罗特方程。现在不管怎么打印输出，一心只关注于实现方程即可。循环中，对\texttt{z}进行平方，然后加上\texttt{c}，知道\texttt{abs}的值小于2。对于一些坐标来说，其值永远不可能比2小，所以当循环次数达到\texttt{max\_iterations}时，我们就决定放弃。最后，将会返回那些\texttt{abs}值收敛的迭代次数：

\begin{codebox}[cpp]
static auto mandelbrot_iterations(cmplx c)
{
    cmplx z {};
    size_t iterations {0};
    const size_t max_iterations {1000};
    while (abs(z) < 2 && iterations < max_iterations) {
        ++iterations;
        z = pow(z, 2) + c;
    }
    return iterations;
}
\end{codebox}

\item 那么现在我们就来实现主函数。在主函数中我们会定义缩放函数对象\texttt{scale}，用于对坐标值进行多维变换：

\begin{codebox}[cpp]
int main()
{
    const size_t w {100};
    const size_t h {40};
    auto scale (scaled_cmplx(
        scaler(0, w, -2.0, 1.0),
        scaler(0, h, -1.0, 1.0)
    ));
\end{codebox}

\item 为了可以在一维上迭代器整个图形，需要完成另一个转换函数，用于将二维图像进行降维操作。其会根据我们所设置的字符宽度进行计算。其会将一维上的长度进行折断，然后进行多行显示，通过使用\texttt{scale}函数对坐标进行变换，然后返回复数坐标：

\begin{codebox}[cpp]
    auto i_to_xy ([=](int i) { return scale(i % w, i / w); });
\end{codebox}

\item 我们将图像的二维坐标(int，int类型)转换为一维坐标(int类型)，再将坐标转换成曼德尔布罗特结合坐标(cmplx类型)。让我们将所有功能放入一个函数，我们将使用一组调用链：

\begin{codebox}[cpp]
    auto to_iteration_count ([=](int i) {
        return mandelbrot_iterations(i_to_xy(i));
    });
\end{codebox}

\item 现在我们可以来设置所有数据。假设我们的结果ASCII图像的字符宽度为\texttt{w}，高度为\texttt{h}。这样就能将结果存储在一个长度为\texttt{w * h}数组中。我们使用\texttt{std::iota}将数值范围进行填充。这些数字可以用来作为转换的输入源 ，我们将变换过程包装在\texttt{to\_iteration\_count}中：

\begin{codebox}[cpp]
vector<int> v (w * h);
iota(begin(v), end(v), 0);
transform(begin(v), end(v), begin(v), to_iteration_count);
\end{codebox}

\item 现在有一个v数组，其使用一维坐标进行初始化，不过后来会被曼德尔布罗特迭代计数所覆盖。因此，我们就可以对图像进行打印。可以将终端窗口设置为\texttt{w}个字符宽度，这样我们就不需要打印换行符。不过，可能会有对\texttt{std::accumulate}有一种创造性的误用。\texttt{std::accumulate}使用二元函数对处理范围进行缩小。我们可以对其提供一个二元函数，其能接受一个输出迭代器(并且我们将在下一步进行终端打印)，并使用范围内的单个值进行计算。如果相应值的迭代次数大于50次时，我们会打印\texttt{*}字符到屏幕上。否则，会打印空字符在屏幕上。在每行结束时(因为计数器变量n可被W均匀地分割)，我们会打印一个换行符：

\begin{codebox}[cpp]
auto binfunc ([w, n{0}] (auto output_it, int x) mutable {
    *++output_it = (x > 50 ? '*' : ' ');
    if (++n % w == 0) { ++output_it = '\n'; }
    return output_it;
});
\end{codebox}

\item 通过对输入范围使用\texttt{std::accumulate}，我们将二元打印函数和\texttt{ostream\_iterator}相结合，我们需要在屏幕上刷新计算出的曼德尔布罗特集合：

\begin{codebox}[cpp]
    accumulate(begin(v), end(v), ostream_iterator<char>{cout},
              binfunc);
}
\end{codebox}

\item 编译并运行程序，就可以看到如下的输出，其看起来和墙上的海报很像吧！

\end{enumerate}

\myGraphic{0.9}{images/chapter6/6-5-3.png}{}

\mySmallsection{How it works...}

整个计算过程都使用\texttt{std::transform}对一维数组进行处理：

\begin{codebox}[cpp]
vector<int> v (w * h);
iota(begin(v), end(v), 0);
transform(begin(v), end(v), begin(v), to_iteration_count);
\end{codebox}

所以，会发生什么呢？我们为什么要这么做？\texttt{to\_iteration\_count}函数是基于从\texttt{i\_to\_xy}开始的调用链，从\texttt{scale}到\texttt{mandelbrot\_iterations}。下面的图像就能展示我们的转换步骤：

\myGraphic{0.9}{images/chapter6/6-5-4.png}{}

这样，我们就可以使用一维数组作为输入，并且获得曼德尔布罗特方程的迭代次数(使用一维坐标表示的二维坐标上的值)。三个互不相关的转换是件好事。这样代码就可以独立的进行测试，这样就不用互相牵制了。同样，这样更容易进行正确性测试，并寻找并修复bug。


% --- recipe 6: chapter6-6-chinese.md ---
\section{实现分割算法}

很多情况下，STL中的算法并不够我们使用，有些算法需要我们自己去实现。解决具体问题之前，我们需要确定，这个问题是否有通解。当我们自己遇到一些问题时，我们可以实现一些辅助函数帮助我们，这些辅助函数逐渐的就可以形成库。这里关键是要明白什么样的代码是足够通用的，否则我们就需要创造一套通用语言了。

本节我们将实现一个算法，叫做\textbf{分割}(split)。该算法可以通过给定的值，对任何范围的元素进行分割，将分割后的结果块拷贝到输出区域中。

\mySmallsection{How to do it...}

本节，将实现类似于STL的算法叫做分割，并且用这个算法对字符串进行分割：

\begin{enumerate}
\item 首先，包含必要的头文件，并声明相应的命名空间。

\begin{codebox}[cpp]
#include <iostream>
#include <string>
#include <algorithm>
#include <iterator>
#include <list>
using namespace std; 
\end{codebox}

\item 本节的所有算法都围绕分割来进行。其接受一对\texttt{begin/end}迭代器和一个输出迭代器，其用法和\texttt{std::copy}或\texttt{std::transform}类似。其他参数为\texttt{split\_val}和\texttt{bin\_func}。\texttt{split\_val}参数是要在输入范围内要查找的值，其表示要当碰到这个值时，要对范围进行分割。\texttt{bin\_func}参数是一个函数，其为分割的子序列的开始和结尾。我们可以使用\texttt{std::find}对输入范围进行迭代查找，这样就能直接跳转到\texttt{split\_val}所在的位置。当将一个长字符串分割成多个单词，可以通过分割空格字符达到目的。对于每一个分割值，都会做相应的分割，并将对应的分割块拷贝到输出范围内：

\begin{codebox}[cpp]
template <typename InIt, typename OutIt, typename T, typename F>
InIt split(InIt it, InIt end_it, OutIt out_it, T split_val,
          F bin_func)
{
    while (it != end_it) {
        auto slice_end (find(it, end_it, split_val));
        *out_it++ = bin_func(it, slice_end);
        if (slice_end == end_it) { return end_it; }
        it = next(slice_end);
    }
    return it;
}
\end{codebox}

\item 现在尝试一下我们的新算法，构建一个需要进行分割的字符串。其中的字符使用\texttt{-}进行连接：

\begin{codebox}[cpp]
int main()
{
    const string s {"a-b-c-d-e-f-g"};
\end{codebox}

\item 创建一个\texttt{bin\_func}对象，其能接受一组迭代器，我们需要通过该函数创建一个新的字符串：

\begin{codebox}[cpp]
    auto binfunc ([](auto it_a, auto it_b) {
        return string(it_a, it_b);
    });
\end{codebox}

\item 输出的子序列将保存在\texttt{std::list}中。我们现在可以调用\texttt{split}算法：

\begin{codebox}[cpp]
list<string> l;
split(begin(s), end(s), back_inserter(l), '-', binfunc);
\end{codebox}

\item 为了看一下结果，我们将对子字符串进行打印：

\begin{codebox}[cpp]
    copy(begin(l), end(l), ostream_iterator<string>{cout, "\n"});
} 
\end{codebox}

\item 编译并运行程序，就可以看到如下输出。其子序列将不会包含破折号，只有单个单词(在我们的例子中，为单个字母)：

\begin{codebox}[cpp]
$ ./split
a
b
c
d
e
f
g
\end{codebox}

\end{enumerate}

\mySmallsection{How it works...}

\texttt{split}算法与\texttt{std::transform}的工作原理很类似，因为其能接受一对\texttt{begin/end}迭代器和一个输出迭代器。其也会将最终的算法结果拷贝到输出迭代器所在的容器。除此之外，其接受一个\texttt{split\_val}值和一个二元函数。让我们再来看一起其整体实现：

\begin{codebox}[cpp]
template <typename InIt, typename OutIt, typename T, typename F>
InIt split(InIt it, InIt end_it, OutIt out_it, T split_val, F bin_func)
{
    while (it != end_it) {
        auto slice_end (find(it, end_it, split_val));
        *out_it++ = bin_func(it, slice_end);
        
        if (slice_end == end_it) { return end_it; }
        it = next(slice_end);
    }
    return it;
}
\end{codebox}

实现中的循环会一直进行到输入范围结束。每次迭代中都会调用\texttt{std::find}用来在输入范围内查找下一个与\texttt{split\_val}匹配的元素。在我们的例子中，分割字符就是\texttt{-}。每次的下一个减号字符的位置会存在\texttt{slice\_end}。每次循环迭代之后，\texttt{it}迭代器将会更新到下一个分割字符所在的位置。循环起始范围将从一个减号跳到下一个减号，而非每一个独立的元素。

这一系列的操作中，迭代器\texttt{it}指向的是最后子字符串的起始位置，\texttt{slice\_end}指向的是子字符串的末尾位置。通过这两个迭代器，就能表示分割后的子字符串。对于字符串\texttt{foo-bar-baz}来说，循环中就有三个迭代器。对于用户而言，迭代器什么的并不重要，他们想要的是子字符串，所以这里就是\texttt{bin\_func}来完成这个任务。当我们调用\texttt{split}时，我们可以给定其一个如下的二元函数：

\begin{codebox}[cpp]
[](auto it_a, auto it_b) {
    return string(it_a, it_b);
}
\end{codebox}

\texttt{split}函数会将迭代器传递给\texttt{bin\_func}，并通过迭代器将结果放入输出迭代器中。这样我们就能通过\texttt{bin\_func}获得相应的单词，这里的结果是\texttt{foo}，\texttt{bar}和\texttt{baz}。

\mySmallsection{There's more...}

我们也可以实现相应的迭代器来完成这个算法的实现。我们现在不会去实现这样一个迭代器，但是可以简单的看一下。

迭代的每次增长，都会跳转到下一个限定符。

当对迭代器进行解引用时，其会通过迭代器指向的当前位置，创建一个字符串对象，就如同之前用到的\texttt{bin\_func}函数那样。

迭代器类可以称为\texttt{split\_iterator}，用来替代算法\texttt{split}，用户的代码可以写成如下的样式：

\begin{codebox}[cpp]
string s {"a-b-c-d-e-f-g"};
list<string> l;

auto binfunc ([](auto it_a, auto it_b) {
    return string(it_a, it_b);
});

copy(split_iterator{begin(s), end(s), "-", binfunc},{}, back_inserter(l));
\end{codebox}

虽然在使用中很方便，但是在实现时，迭代器的方式要比算法的形式复杂许多。并且，迭代器实现中很多边缘值会触发代码的bug，并且迭代器实现需要经过非常庞杂的测试。不过，其与其他STL算法能够很好的兼容。


% --- recipe 7: chapter6-7-chinese.md ---
\section{将标准算法进行组合}

\texttt{gather}为STL算法中最好的组合性例子。Sean Parent在任Adobe系统首席科学家时，就在向世人普及这个算法，因为其本身短小精悍。其使用方式就如同做一件艺术品一样。

\texttt{gather}算法能操作任意的元素类型。其更改元素的顺序，通过用户的选择，其会将对应的数据放置在对应位置上。

\mySmallsection{How to do it...}

本节，我们来实现\texttt{gather}算法，并对其进行一定的修改。最后，将展示如何进行使用：

\begin{enumerate}
\item 包含必要的头文件，声明所使用的命名空间。

\begin{codebox}[cpp]
#include <iostream>
#include <algorithm>
#include <string>
#include <functional>
using namespace std; 
\end{codebox}

\item \texttt{gather}算法是表现标准算法组合性很好的一个例子。\texttt{gather}接受一对\texttt{begin/end}迭代器和另一个迭代器\texttt{gather\_pos}，其指向\texttt{begin}和\texttt{end}迭代器的中间位置。最后一个参数是一个谓词函数。使用谓词函数，算法会将满足谓词条件的所有元素放置在\texttt{gather\_pos}迭代器附近。使用\texttt{std::stable\_partition}来完成移动元素的任务。\texttt{gather}算法将会返回一对迭代器。这两个迭代器由\texttt{stable\_partition}所返回，其表示汇集范围的起始点和结束点：

\begin{codebox}[cpp]
template <typename It, typename F>
pair<It, It> gather(It first, It last, It gather_pos, F predicate)
{
    return {stable_partition(first, gather_pos, not_fn(predicate)),
            stable_partition(gather_pos, last, predicate)};
}
\end{codebox}

\item 算法的另一种变体为\texttt{gather\_sort}。其工作原理与\texttt{gather}相同，不过其不接受一元谓词函数；其能接受一个二元比较函数。这样，其就也能将对应值汇集在\texttt{gather\_pos}附近，并且能知道其中最大值和最小值：

\begin{codebox}[cpp]
template <typename It, typename F>
void gather_sort(It first, It last, It gather_pos, F comp_func)
{
    auto inv_comp_func ([&](const auto &...ps) {
        return !comp_func(ps...);
    });
    stable_sort(first, gather_pos, inv_comp_func);
    stable_sort(gather_pos, last, comp_func);
}
\end{codebox}

\item 让我们来使用一下这些算法。先创建一个谓词函数，其会告诉我们当前的字母是不是\texttt{a}，再构建一个字符串，仅包含\texttt{a}和\texttt{-}字符：

\begin{codebox}[cpp]
int main()
{
    auto is_a ([](char c) { return c == 'a'; });
    string a {"a_a_a_a_a_a_a_a_a_a_a"};
\end{codebox}

\item 继续构造一个迭代器，让其指向字符串的中间位置。这时可以调用\texttt{gather}算法，然后看看会发生什么。\texttt{a}字符将汇集在字符串中间的某个位置附近：

\begin{codebox}[cpp]
    auto middle (begin(a) + a.size() / 2);
    gather(begin(a), end(a), middle, is_a);
    cout << a << '\n';
\end{codebox}

\item 再次调用\texttt{gather}，不过这次\texttt{gather\_pos}的位置在字符串的起始端：

\begin{codebox}[cpp]
    gather(begin(a), end(a), begin(a), is_a);
    cout << a << '\n';
\end{codebox}

\item 再将\texttt{gather\_pos}的位置放在末尾试试：

\begin{codebox}[cpp]
    gather(begin(a), end(a), end(a), is_a);
    cout << a << '\n';
\end{codebox}

\item 最后一次，这次将再次将迭代器指向中间位置。这次与我们的期望不相符，后面我们来看看发生了什么：

\begin{codebox}[cpp]
    // This will NOT work as naively expected
    gather(begin(a), end(a), middle, is_a);
    cout << a << '\n';
\end{codebox}

\item 再构造另一个字符串，使用下划线和一些数字组成。对于这个输入队列，我们使用\texttt{gather\_sort}。\texttt{gather\_pos}迭代器指向字符串的中间，并且比较函数为\texttt{std::less<char>}：

\begin{codebox}[cpp]
    string b {"_9_2_4_7_3_8_1_6_5_0_"};
    gather_sort(begin(b), end(b), begin(b) + b.size() / 2,
               less<char>{});
    cout << b << '\n';
}
\end{codebox}

\item 编译并运行程序，就会看到如下的输出。对于前三行来说和预期一样，不过第四行貌似出现了一些问题。最后一行中，我们可以看到\texttt{gather\_short}函数的结果。数字的顺序是排过序的(中间小，两边大)：

\begin{codebox}[cpp]
$ ./gather
_____aaaaaaaaaaa_____
aaaaaaaaaaa__________
__________aaaaaaaaaaa
__________aaaaaaaaaaa
_____9743201568______
\end{codebox}

\end{enumerate}

\mySmallsection{How it works...}

\texttt{gather}算法有点难以掌握，因为其非常短，但处理的是比较复杂的问题。让我们来逐步解析这个算法：

\myGraphic{0.9}{images/chapter6/6-7-1.png}{}

\begin{enumerate}
\item 初始化相应元素，并提供一个谓词函数。图中满足谓词函数条件的元素为灰色，其他的为白色。迭代器\texttt{a}和\texttt{c}表示整个范围的长度，并且迭代器\texttt{b}指向了最中间的元素。这就表示要将所有灰色的格子聚集在这个迭代器附近。

\item \texttt{gather}算法会对范围\texttt{(a, b]}调用\texttt{std::stable\_partition}，并对另一边使用不满足谓词条件的结果。这样就能将所有灰色格子集中在b迭代器附近。

\item 另一个\texttt{std::stable\_partition}已经完成，不过在\texttt{[b, c)}间我们将使用满足谓词函数的结果。这样就将灰色的格子汇集在b迭代器附近。

\item 所有灰色的格子都汇集在\texttt{b}迭代器附近，这是就可以返回起始迭代器\texttt{x}和末尾迭代器\texttt{y}，用来表示所有连续灰色格子的范围。

\end{enumerate}

我们对同一个范围多次调用\texttt{gather}算法。最初，将所有元素放在范围中间位置。然后，尝试放在开始和末尾。这种实验很有趣，因为这会让其中一个\texttt{std::stable\_partition}无元素可以处理。

最后一次对\texttt{gather}进行调用时，参数为\texttt{(begin, end, middle)}，但没有和我们预期的一样，这是为什么呢？这看起来像是一个bug，实际上不是。

试想一个字符组\texttt{aabb}，使用谓词函数\texttt{is\_character\_a}，用来判断元素是否为\texttt{a}，当我们将第三个迭代器指向字符范围的中间时，会复现这个bug。原因：第一个\texttt{stable\_partition}调用会对子范围\texttt{aa}进行操作，并且另一个\texttt{stable\_partition}会对子范围\texttt{bb}上进行操作。这种串行的调用时无法得到\texttt{baab}的，其结果看起来和开始一样，没有任何变化。

\begin{myInfoBox}
Note： 要是想得到我们预期的序列，我们可以使用\texttt{std::rotate(begin, begin + 1, end);}
\end{myInfoBox}

\texttt{gather\_sort}基本上和\texttt{gather}差不多。签名不同的就是在于谓词函数。实现的不同在于\texttt{gather}调用了两次\texttt{std::stable\_partition}，而\texttt{gather\_sort}调用了两次\texttt{std::stable\_sort}。

这是由于\texttt{not\_fn}不能作用域二元函数，所以反向比较不能由\texttt{not\_fn}完成。


% --- recipe 8: chapter6-8-chinese.md ---
\section{删除词组间连续的空格}

我们会经常从输入中读取字符串，这些字符串会包含一些原生格式，需要进行清洗。其中一个例子就是字符串中包含了太多的空格。

本节，我们将实现一个聪明的空格滤波算法，其会删除多于的空格，会给单词间留下一个空格。我们可以将这个算法称为\texttt{remove\_multi\_whitespace}，并且接口与STL很像。

\mySmallsection{How to do it...}

本节，我们将实现过滤空格的算法，并了解其是如何进行工作的：

\begin{enumerate}
\item 包含必要的头文件和声明所使用的命名空间：

\begin{codebox}[cpp]
#include <iostream>
#include <string>
#include <algorithm>
using namespace std;
\end{codebox}

\item \texttt{remove\_multi\_whitespace}看起来与STL的风格非常类似。这个算法会移除多余的空格，只保留一个空格。当字符串为\texttt{a b}，算法是不会进行任何操作的；当字符串为\texttt{a b}时，算法会返回\texttt{a b}。为了完成这个算法，我们使用\texttt{std::unqiue}通过对一段区域的迭代，用来查找一对连续的元素。然后，通过谓词函数进行判断，确定两个元素是否相等。如果相等，那么\texttt{std::unique}会将其中一个移除。这样，子范围中就不会存在相等的元素了。谓词函数会通过读取到的内容来判断二者是否相等。我们需要给\texttt{std::unique}怎么样一个谓词函数呢？其需要判断两个元素是否是连续的空格；如果是，就要移除一个空格。与\texttt{std::unique}类似，也需要传入一对\texttt{begin/end}迭代器，然后返回的迭代器将返回新范围的末尾迭代器：

\begin{codebox}[cpp]
template <typename It>
It remove_multi_whitespace(It it, It end_it)
{
    return unique(it, end_it, [](const auto &a, const auto &b) {
        return isspace(a) && isspace(b);
    });
}
\end{codebox}

\item 万事俱备，就来进行测试，尝试使用算法将不必要的空格进行删除：

\begin{codebox}[cpp]
int main()
{
    string s {"fooo bar \t baz"};
    cout << s << '\n';
\end{codebox}

\item 对字符串使用过滤算法，去掉多余的空格：

\begin{codebox}[cpp]
    s.erase(remove_multi_whitespace(begin(s), end(s)), end(s));
    cout << s << '\n';
}
\end{codebox}

\item 编译并运行程序，就会得到如下的输出：

\begin{codebox}[cpp]
$ ./remove_consecutive_whitespace
fooo bar       baz
fooo bar baz
\end{codebox}

\end{enumerate}

\mySmallsection{How it works...}

整个问题的解决中，我们没有使用循环或者元素间的互相比较。我们只使用谓词函数来完成判断两个给定字符是否是空格的任务。然后，将谓词函数与\texttt{std::unique}相结合，所有多余的空格就都消失了。本章中有些算法可能会有些争议，不过这个算法的确算的上短小精悍的典范了。

我们如何在将算法进行组合的呢？我们来看一下\texttt{std::unique}可能的实现代码：

\begin{codebox}[cpp]
template<typename It, typename P>
It unique(It it, It end, P p)
{
    if (it == end) { return end; }

    It result {it};
    while (++it != end) {
        if (!p(*result, *it) && ++result != it) {
            *result = std::move(*it);
        }
    }
    return ++result;
}
\end{codebox}

其中循环会迭代到范围的最后，当元素满足谓词条件，就会从原始位置上移除一个元素。这个版本的\texttt{std::unique}不接受多余的谓词函数，来判断两个相邻的元素是否相等。这样的话，只能将重复的字符去除，比如会将\texttt{abbbbbbbbc}变换成\texttt{abc}。

那么，我们应该怎么做才能不去除除了空格之外的重复的元素呢？这样，谓词函数不能告诉程序“两个输入字符是相同的”，而是要说“两个输入字符都是空格”。

最后需要注意的是，无论是\texttt{std::unique}还是\texttt{remove\_multi\_whitespace}都会从字符串中移除字母元素。根据字符串的语义来移动字符串，并表明新的结尾在哪里。新的尾部到旧的尾部的元素依旧存在，所以我们必须将它们删除：

\begin{codebox}[cpp]
s.erase(remove_multi_whitespace(begin(s), end(s)), end(s));
\end{codebox}

和\texttt{vector}和\texttt{list}一样，\texttt{erase}成员函数其会对元素进行擦除和删除。


% --- recipe 9: chapter6-9-chinese.md ---
\section{压缩和解压缩字符串}

压缩问题在编程面试中出现的相对较多。就是使用一个函数将\texttt{aaaaabbbbbbbccc}字符串转换成一个短字符串\texttt{a5b7c3}。\texttt{a5}表示原始字符串中有5个a，\texttt{b7}表示原始字符串中有7个b。这就一个相对简单的压缩算法。对于普通的文本，并不需要使用这个算法，因为文本中重复的东西很少，不需要进行压缩。不过，这套算法就算没有计算机，我们也能很容易的对其进行实现。如果代码在一开始没有进行很好的设计，那么就很容易出现bug。虽然，处理字符串并不是一件很困难的事情，但是代码中大量使用C风格的字符串时，很有可能遇到缓冲区溢出的bug。

本节让我们使用STL来对字符压缩和解压进行实现。

\mySmallsection{How to do it...}

本节，我们将对字符串实现简单的\texttt{compress}和\texttt{decompress}函数：

\begin{enumerate}
\item 包含必要的头文件，并声明所使用的命名空间：

\begin{codebox}[cpp]
#include <iostream>
#include <string>
#include <algorithm>
#include <sstream>
#include <tuple>
using namespace std;
\end{codebox}

\item 对于我们的压缩算法，我们会尝试去找到文本中连续相同的字符，并且对他们进行单独的进行压缩处理。当我们拿到一个字符串，我们需要知道与第一个字符不同的字符在哪里。这里使用\texttt{std::find}来寻找与第一个位置上的元素不同的元素位置。先将起始位置的字符赋予\texttt{c}。经过查找后就会返回一个迭代器，其指向第一个不同的元素。两个不同字符间的距离，会放到元组中返回：

\begin{codebox}[cpp]
template <typename It>
tuple<It, char, size_t> occurrences(It it, It end_it)
{
    if (it == end_it) { return {it, '?', 0}; }
    const char c {*it};
    const auto diff (find_if(it, end_it,
                    [c](char x) { return c != x; }));
    return {diff, c, distance(it, diff)};
}
\end{codebox}

\item \texttt{compress}会连续的对\texttt{occurrences}函数进行调用。这样，就能从同一个字符组，跳转到另一个。\texttt{r << c << n}行表示将字符\texttt{c}推入到输出流中，并且将\texttt{occurrences}函数的调用次数作为结果字符串的一部分。最后会返回一个字符串对象，就包含了压缩过的字符串：

\begin{codebox}[cpp]
string compress(const string &s)
{
    const auto end_it (end(s));
    stringstream r;
    for (auto it (begin(s)); it != end_it;) {
        const auto [next_diff, c, n] (occurrences(it, end_it));
        r << c << n;
        it = next_diff;
    }
    return r.str();
}
\end{codebox}

\item \texttt{decompress}的原理也不复杂，但会更简短。其会持续的从输入流中获取字符，字符串包括字符和数字。对于这两种值，函数会构造一个字符串用于解压所获取到的字符串。最后，会再次返回一个字符串。顺带一提，这里的\texttt{decompress}函数是不安全的。其很容易被破解。我们会在后面来看下这个问题：

\begin{codebox}[cpp]
string decompress(const string &s)
{
    stringstream ss{s};
    stringstream r;
    char c;
    size_t n;
    while (ss >> c >> n) { r << string(n, c); }
    return r.str();
}
\end{codebox}

\item 主函数中会构造一个简单的字符串，里面有很多重复的字符。打印压缩过后，和解压过后的字符串。最后，我们应该会得到原始的字符串：

\begin{codebox}[cpp]
int main()
{
    string s {"aaaaaaaaabbbbbbbbbccccccccccc"};
    cout << compress(s) << '\n';
    cout << decompress(compress(s)) << '\n';
}
\end{codebox}

\item 编译并运行程序，我们就会得到如下的输出：

\begin{codebox}[cpp]
$ ./compress
a9b9c11
aaaaaaaaabbbbbbbbbccccccccccc
\end{codebox}

\end{enumerate}

\mySmallsection{How it works...}

这里我们使用两个函数\texttt{compress}和\texttt{decompress}来解决这个问题。

解压函数这里实现的十分简单，因为其就包含一些变量的声明，其主要工作的代码其实只有一行：

\begin{codebox}[cpp]
while (ss >> c >> n) { r << string(n, c); }
\end{codebox}

其能持续将字符读取到\texttt{c}当中，并且将数字变量读取到\texttt{n}中，然后输出到\texttt{r}中。\texttt{stringstream}类在这里会隐藏对字符串解析的细节。当成功进行解压后，解压的字符串将输入到字符流中，这也就是\texttt{decompress}最后的结果。如果\texttt{c = 'a'}并且\texttt{n = 5}，那么\texttt{string(n, c)}的字符串为\texttt{aaaaa}。

压缩函数比较复杂，我们为其编写了一个小的辅助函数。这个辅助函数就是\texttt{occurences}。那么我们就先来看一下\texttt{occurences}函数。下面的图展示了\texttt{occurences}函数工作的方式：

\myGraphic{0.9}{images/chapter6/6-9-1.png}{}

\texttt{occurences}函数能够接受两个参数：指向字符序列起始点的迭代器和末尾点的迭代器。使用\texttt{find\_if}能找到第一个与起始点字符不同的字符的位置，也就是图中的\texttt{diff}迭代器的位置。起始位置与\texttt{diff}位置之间元素就与起始字符相同，图中相同的字符有6个。在我们计算出这些信息后，\texttt{diff}迭代就可以在下次查询时，进行重复利用。因此，我们将\texttt{diff}、子序列范围和子序列范围的长度包装在一个元组中进行返回。

根据这些信息，我们就能在子序列之间切换，并且将结果推入到目标字符串中：

\begin{codebox}[cpp]
for (auto it (begin(s)); it != end_it;) {
    const auto [next_diff, c, n] (occurrences(it, end_it));
    r << c << n;
    it = next_diff;
}
\end{codebox}

\mySmallsection{There's more...}

还记得在第4步的时候，我们说过\texttt{decompress}不安全吗？这个函数确实容易被利用。

试想我们传入一个字符串：\texttt{a00000}。压缩的第一个结果为\texttt{a1}因为其只包含了一个字母\texttt{a}。然后，对后面5个0进行处理，结果为\texttt{05}。然后将两个结果合并，那么结果就为\texttt{a105}。不幸的是，外部对这个字符串的解读是“a连续出现了105次”。我们的输入字符串并没有什么错。这里最糟糕的情况就是，我们将这个字符串进行了压缩，然后我们通过输入的六个字符得到了一个长度为105的字符串。试想当用户得到了这样的结果会不会感到愤怒？因为我们的算法并没有准备好应对这样的输入。

为了避免这样的事情发生，我们只能在\texttt{compress}函数中禁止数字的输入，或者将数字使用其他的方式进行处理。之后，\texttt{decompress}算法需要加入一个条件，就是需要固定输出字符串的最大长度。这个就当做作业，交由读者自行完成。


